% ============================================================
% PHẦN: KIẾN TRÚC SERVER VÀ GAME  (4 slide)
% Nội dung lấy từ file "Kiến Trúc Server & Hướng Dẫn Của Đức.pptx" (Thành viên A: Đức),
% đã ĐỐI CHIẾU với code thật trên main (TcpServer.cs, ClientSession.cs, RoomManager.cs, RoomHandler.cs).
% Chỉ ghi những gì đã hiện thực; phần chưa xong ghi rõ "đang phát triển".
% ============================================================
\section{Kiến Trúc Server Và Game}

% ---------- Slide: lõi TCP ----------
\begin{frame}{Lõi TCP Và Quản Lý Phiên}
  \begin{columns}[T]
    \begin{column}{0.52\textwidth}
      \begin{block}{Cách server xử lý kết nối}
        \begin{itemize}
          \item \texttt{TcpListener} lắng nghe bất đồng bộ, \textbf{mỗi client một Task} riêng.
          \item \textbf{Framing 4-byte}: đọc đủ \texttt{PayloadLength} rồi mới giải mã gói.
          \item \textbf{Khóa luồng gửi} \texttt{SemaphoreSlim(1,1)}: chặn hai Task ghi xen kẽ byte vào cùng một stream.
          \item \texttt{PacketRouter} chuyển gói tới từng \texttt{IPacketHandler} đã đăng ký.
        \end{itemize}
      \end{block}
    \end{column}
    \begin{column}{0.44\textwidth}
      \centering
      \begin{tikzpicture}[
        flow/.style={draw=inkblue, fill=white, rounded corners=3pt, thick, align=center,
                     minimum height=0.85cm, minimum width=5.2cm, font=\scriptsize},
        arr/.style={-{Stealth[length=2mm]}, thick, color=brandblue}]
        \node[flow, fill=paleblue] (a) at (0,0)     {1. \texttt{NetworkStream} $\rightarrow$ \texttt{PacketFraming} (4-byte)};
        \node[flow] (b) at (0,-1.45) {2. \texttt{ClientSession.OnPacketReceived}};
        \node[flow] (c) at (0,-2.9)  {3. \texttt{PacketRouter.RouteAsync(session, pkt)}};
        \node[flow, fill=paleblue] (d) at (0,-4.35) {4. \texttt{RoomHandler.HandleAsync(session, pkt)}};
        \draw[arr] (a) -- (b);
        \draw[arr] (b) -- (c);
        \draw[arr] (c) -- (d);
      \end{tikzpicture}

      \vspace{0.1cm}
      \scriptsize Luồng đi của một gói tin tại server.
    \end{column}
  \end{columns}
\end{frame}

% ---------- Slide: quản lý phòng ----------
\begin{frame}{Quản Lý Phòng Đấu}
  \begin{columns}[T]
    \begin{column}{0.50\textwidth}
      \begin{block}{\texttt{RoomManager}}
        \begin{itemize}
          \item \texttt{ConcurrentDictionary}: lưu danh sách phòng an toàn đa luồng.
          \item \textbf{Khóa riêng từng phòng}: hai người vào cùng một lúc không bị lỗi.
          \item \textbf{Host thoát} thì phòng bị giải tán và người còn lại được báo (\texttt{RoomUpdate}).
        \end{itemize}
      \end{block}
    \end{column}
    \begin{column}{0.46\textwidth}
      \footnotesize
      \renewcommand{\arraystretch}{1.25}
      \begin{tabularx}{\textwidth}{@{}Y Y@{}}
        \toprule
        \textbf{Gói tin} & \textbf{Chức năng} \\
        \midrule
        \texttt{GetRoomListRequest} & Danh sách phòng \\
        \texttt{CreateRoomRequest} & Tạo phòng bằng ID \\
        \texttt{JoinRoomRequest} & Vào phòng \\
        \texttt{LeaveRoom} & Rời phòng \\
        \bottomrule
      \end{tabularx}
      \vspace{0.15cm}

      \scriptsize \textbf{Đang phát triển:} nối chức năng Kick (đã có ở \texttt{RoomManager}) vào gói tin.
    \end{column}
  \end{columns}
\end{frame}

% ---------- Slide: state machine ----------
\begin{frame}{Máy Trạng Thái Của Phòng Đấu}
  \centering
  \vspace{0.6cm}
  \slideimg[0.97\textwidth]{12-state-machine.png}{Sơ đồ State Machine}

  \vspace{0.4cm}
  \footnotesize Hai trạng thái đầu \textbf{đã hiện thực}; ba trạng thái sau đang phát triển.
\end{frame}

% ---------- Slide: bảo mật + ngắt kết nối ----------
\begin{frame}{Bảo Mật Và Xử Lý Ngắt Kết Nối}
  \begin{columns}[T]
    \begin{column}{0.49\textwidth}
      \begin{block}{Dọn dẹp khi đứt mạng}
        \begin{itemize}
          \item Bắt \texttt{IOException} và \texttt{SocketException} \textbf{riêng biệt}.
          \item Khối \texttt{finally} luôn gọi \texttt{Cleanup()} để đóng luồng mạng.
          \item Sự kiện \texttt{OnDisconnected}: xóa khỏi danh sách online và \textbf{báo Offline} cho mọi người.
        \end{itemize}
      \end{block}
    \end{column}
    \begin{column}{0.47\textwidth}
      \begin{exampleblock}{Mã hóa đường truyền}
        \begin{itemize}
          \item Toàn bộ payload JSON qua TCP được mã hóa \textbf{AES-128}.
          \item Gói chữ rõ không có khóa bị server \textbf{đóng kết nối}.
        \end{itemize}
      \end{exampleblock}
      \begin{alertblock}{Đang phát triển}
        \footnotesize Tự động cho người chơi rời phòng khi mất mạng giữa chừng.
      \end{alertblock}
    \end{column}
  \end{columns}
\end{frame}
